% ESTILO: IEEE (biblatex-ieee)
% cita-ieee-config.tex
\usepackage[
  style=ieee,
  sortcites=true,
  citereset=chapter,
  backend=biber,
  backref=false,
  indexing=true
]{biblatex}

% ============================================================
% LAS DOS CADENAS PROPIAS DE biblatex-ieee
% ------------------------------------------------------------
% ieee.bbx LAS DECLARA (\NewBibliographyString) Y LAS DEFINE SOLO
% EN INGLÉS (v1.3f, RF-14). EN CASTELLANO biblatex IMPRIME LA CLAVE
% EN NEGRITA («presentedat», MEDIDO EN EL PDF DE DUEK, CON EL AVISO
% «Bibliography string 'presentedat' untranslated» EN EL .log), Y EN
% PANTALLA ESA NEGRITA CORTA LA HOJA (SC-57).
% presentedat: DECISIÓN DE ALBERTO. patentjp: LA FORMA DE LAS
% DEMÁS PATENTES DE LA EDITORIAL (cita-verona-modern-config.tex).
% LAS CADENAS COMUNES A TODOS LOS ESTILOS ESTÁN EN cita-cadenas-es.tex.
% ============================================================
\DefineBibliographyStrings{spanish}{%
  presentedat = {presentado en\addspace},
  patentjp    = {patente japonesa},
}

% ============================================================
% LA COMA Y EL PUNTO DESPUÉS DE LA COMILLA DE CIERRE
% ------------------------------------------------------------
% ieee.bbx LLAMA A \uspunctuation SIN DISTINGUIR IDIOMA (v1.3f,
% LÍNEA 34): LA CONVENCIÓN NORTEAMERICANA, QUE METE LA COMA Y EL
% PUNTO DENTRO DE LAS COMILLAS («título,»). EN CASTELLANO VAN
% DESPUÉS («título»,). \stdpunctuation (biblatex.def) DESHACE LO
% QUE HACE \uspunctuation, Y COMO EXTRA DEL ESPAÑOL SOLO ACTÚA EN
% UN LIBRO EN ESPAÑOL: UNO EN INGLÉS CONSERVA LA CONVENCIÓN DE
% IEEE. MEDIDO CON EL .bib DE DUEK: 87 COMAS Y 5 PUNTOS CAMBIAN DE
% LADO Y NADA MÁS
% ============================================================
\DefineBibliographyExtras{spanish}{\stdpunctuation}
